\subsection{Materials and Structural Models}
\label{subsec:structures}

Three lead-free tin-halide perovskite families were studied:
\CsSnI{}, \FASnI{}, and \MASnI{}.
For each material, three polymorphs were considered: the
low-temperature orthorhombic phase, the intermediate-temperature
tetragonal phase, and the high-temperature cubic phase.
Initial equilibrium structures for the all-inorganic system and for
the orthorhombic \FASnI{} and \MASnI{} phases were taken from the
Materials Project database.\cite{mp2013}
Missing hybrid-phase structures (tetragonal and cubic FA and MA
variants) were constructed by substituting the organic cation into
the corresponding \CsSnI{} prototype and re-relaxing with DFT at
the level described below.
This procedure yielded nine equilibrium phase structures (three
materials $\times$ three polymorphs) as the starting point for all
subsequent calculations.

For each equilibrium structure, three lattice states were prepared:
the equilibrium volume ($\alpha = 1.00$), an isotropic linear
expansion of \SI{3}{\percent} ($\alpha = 1.03$, corresponding to
approximately \SI{9}{\percent} volume expansion), and an isotropic
linear compression of \SI{3}{\percent} ($\alpha = 0.97$).
These 27 starting structures were used for candidate-configuration
generation and model preparation as described in
Section~\ref{subsec:mlworkflow}.

Supercell sizes were chosen so that the minimum lattice dimension
exceeded \SI{25}{\angstrom} in all cases.
Phase-behaviour simulations (Task~A) used large supercells
containing 8{,}232--10{,}368 atoms to support statistically
meaningful structural averaging over long trajectories.
Electronic-structure studies (Tasks~B, C, and D) used smaller
supercells of 625--1{,}536 atoms to enable large-scale SS-DFT
evaluation.
Full supercell dimensions and atom counts are provided in
Table~S1 of the Supporting Information.

\subsection{Electronic-Structure Settings}
\label{subsec:dft_settings}

All DFT and SS-DFT calculations used the following common settings:
\begin{itemize}
  \item Exchange-correlation functional: Perdew--Burke--Ernzerhof
        (PBE) generalised gradient approximation.\cite{pbe1996}
  \item Pseudopotentials: optimised norm-conserving Vanderbilt
        (ONCV) pseudopotentials from the SPMS database.\cite{oncv}
  \item Real-space grid spacing: \SI{0.105}{\angstrom} (equivalent
        to a planewave cutoff of approximately \SI{1500}{eV}).
  \item $k$-point sampling: $\Gamma$-point only, consistent with
        the large supercell sizes used.
  \item Spin--orbit coupling (SOC): not included in the main study
        (see below).
  \item Ionic relaxation convergence: maximum force threshold of
        \SI{0.02}{eV\per\angstrom}.
\end{itemize}

Spin--orbit coupling was not included in the main study.
To bracket the magnitude of this approximation, we performed an
exchange-correlation benchmark on a 384-atom triclinic \MASnI{}
supercell (Materials Project structure mp-995238) using five
functionals.
The results are summarised in Table~S2.
At the PBE level, the band gap is \SI{0.565}{eV}, which lies near
the centre of the GGA range of \SIrange{0.463}{0.646}{eV}
(RevPBE, PBEsol, PBE, RPBE).
Including SOC at the PBE+SOC level reduces the gap substantially to
\SI{0.230}{eV}, consistent with the known strong spin--orbit
effect in tin-containing compounds.\cite{ref_soc_sn}
Accordingly, all results in this work are reported at the PBE level
and should be interpreted as relative trends and mechanism
separation rather than as absolute band-gap values.

\subsection{Machine-Learning Workflow and Model Preparation}
\label{subsec:mlworkflow}

To enable both large-scale finite-temperature MD and efficient
electronic-structure evaluation, we deployed a two-component
machine-learning workflow comprising a system-specific
electronic-structure model for SS-DFT and a task-specific MLFF
for long-time MD.

\subsubsection{Candidate Configuration Generation}

A pool of 1{,}600 candidate configurations was generated by running
short \textit{NVT} molecular dynamics at \SI{400}{K} on each of
the 27 starting structures using the universal force field
MACE-OMAT-medium,\cite{ref_mace_omat} with a \SI{1.5}{fs}
timestep for 4{,}000 steps and one frame sampled every 200 steps.
This generated a chemically diverse pool spanning all relevant
phases, lattice states, and moderate local distortions.

\subsubsection{Representative-Structure Selection and DFT Labelling}

An automatic sampling algorithm was applied to select 84
representative structures from the 1{,}600-structure pool,
minimising redundancy in local chemical-environment space while
maximising coverage of the relevant structural diversity.
Full DFT calculations were performed on these 84 representative
structures to provide reference data for model training.

\subsubsection{SS-DFT Model Training and Validation}

A system-specific electronic-structure model was trained on the
84 DFT-labelled structures.
The trained model predicts the converged electron density without a
full self-consistent field (SCF) iteration, enabling SS-DFT
band-gap evaluation at a fraction of the cost of full DFT.
Validation against held-out full DFT calculations (18 structures
spanning all three materials and both strained and unstrained
conditions) gives a charge-density relative $L_2$ error of 0.0027
and a mean SS-DFT band-gap error of \SI{0.018}{eV}
(Figure~S1 in the Supporting Information), confirming that the
model is suitable for comparative band-gap analysis.

\subsubsection{ML Force Field Training and Validation}

The trained electronic-structure model was used to run ML-DFT
predictions on the full 1{,}600-structure candidate pool.
These energies, forces, and stresses, combined with the 84 full
DFT references, were used to train a task-specific MLFF using the
neuroevolution potential (NEP) architecture.\cite{ref_nep}
A D3 dispersion correction was included for the main production runs
on the hybrid systems, which improved long-time simulation
stability.
Validation against full DFT on the 18-structure held-out set gives:
energy RMSE of \SI{4.79}{eV} per system,
force RMSE of \SI{0.069}{eV\per\angstrom}, and
stress error of \SI{3.5e-4}{eV\per\cubic\angstrom}
(Figures~S2--S4), confirming that the MLFF is suitable for the
finite-temperature MD studies.

\subsection{Phase-Behaviour Study (Task~A)}
\label{subsec:taskA}

Finite-temperature phase behaviour was studied using long-time
\textit{NPT} MD with the trained MLFF.
Heating and cooling cycles ran from \SI{0.1}{K} to \SI{550}{K}
and back in \SI{50}{K} increments.
At each temperature, \SI{50}{ps} of equilibration was followed by
\SI{200}{ps} of production MD at a \SI{0.5}{fs} timestep, with
structural snapshots collected every \SI{10}{ps}.
A stochastic velocity rescaling (SVR) thermostat and stochastic
cell rescaling (SCR) barostat were used, with orthorhombic
cell-shape updates enabled.

Phase transitions were identified from the temperature-dependent
evolution of:
(i) the lattice-parameter ratio $b/a$, which approaches unity as
the structure becomes more tetragonal;
(ii) direction-resolved Sn--I--Sn framework linkage distances along
the three crystallographic axes; and
(iii) Glazer-angle statistics characterising octahedral tilt
magnitude and disorder.

A follow-up calculation for \CsSnI{} was carried out using a
\ce{Cs}-specific refined model without D3 correction, because the
baseline all-inorganic run under the common protocol became unstable
above approximately \SI{400}{K}.
This follow-up is presented separately in the Results section and
does not form part of the common production protocol.

\subsection{Lattice- and Phase-Dependent Band Gap (Task~B)}
\label{subsec:taskB}

Lattice-expansion and phase-dependence effects on band gap were
studied by applying controlled isotropic volume scaling to each of
the nine equilibrium phase structures.
The volume-scaling factor $\alpha$ was varied from 0.90 to 1.10
in steps of 0.025.
For each scaled structure, lattice vectors were fixed, internal
atomic coordinates were relaxed to a force convergence of
\SI{0.02}{eV\per\angstrom} using the MLFF, and the band gap was
evaluated with SS-DFT.

\subsection{Cation-Orientation Sensitivity Study (Task~C)}
\label{subsec:taskC}

Task~C was performed for \FASnI{} and \MASnI{} only,
in two phase regimes.

\subsubsection{Tetragonal-Phase Study}

Cation orientations were varied using an ordered, synchronised
rotation protocol guided by the two-fold mirror symmetry of the
tetragonal structure, as observed experimentally in ordered hybrid
perovskites.\cite{schueller2018crystal}
For each rotation angle, the inorganic cage was relaxed with the
MLFF and the band gap was evaluated with SS-DFT.
An additional randomly rotated configuration was prepared for each
system as a disordered reference.

\subsubsection{High-Temperature Cubic-Phase Study}

A sweep over 20 mean-angle configurations was performed by varying
the mean polar and azimuthal angles of the cation distribution
(polar: \ang{0}, \ang{43}, \ang{131}, \ang{175};
azimuthal: \ang{0}, \ang{88}, \ang{177}, \ang{266}, \ang{355}).
For each configuration, the inorganic cage was relaxed and the band
gap was evaluated with SS-DFT.

\subsection{Integrated Temperature-and-Strain Study (Task~D)}
\label{subsec:taskD}

Task~D combined finite-temperature structural sampling with
externally imposed strain.
For each material and phase, the system was heated from \SI{10}{K}
to the target temperature using an \textit{NVT} simulation with a
Nos\'{e}--Hoover chain (NHC) thermostat (\SI{100}{ps}), followed
by \textit{NPT} equilibration to relax the lattice (SCR barostat,
\SI{100}{ps}).
A configuration was then drawn from the \textit{NPT} trajectory,
and three lattice variants were prepared:
compressive strain ($\alpha = 0.95$), no strain ($\alpha = 1.00$),
and tensile strain ($\alpha = 1.05$).
Each strained structure was re-equilibrated at the target
temperature (\textit{NVT}, \SI{50}{ps}), and two representative
snapshots were extracted and evaluated with SS-DFT.
Representative temperatures were
\SI{250}{K}/\SI{325}{K}/\SI{400}{K}
(orthorhombic/tetragonal/cubic) for \CsSnI{}, and
\SI{150}{K}/\SI{250}{K}/\SI{320}{K} for
\FASnI{} and \MASnI{}.
These temperatures are representative study conditions for each
phase regime; they are not calibrated phase-transition temperatures
within the present workflow.